\chapter{Control Algorithms\label{chap:algorithms}}

This chapter presents the two core control algorithms used in the system: the pursuit control algorithm in \texttt{catch\_executor} and the following control algorithm in \texttt{follower\_manager}.

\section{Pursuit Control Algorithm}

The catch executor uses a \textbf{turn-then-drive} control strategy. At each control tick (default: 20\,Hz), the following computations are performed:

\subsection{Geometric Computations}

Given the master turtle's position $(x_m, y_m, \theta_m)$ and the target's position $(x_t, y_t)$:

\begin{enumerate}
  \item Compute the displacement:
  \[
    \Delta x = x_t - x_m, \quad \Delta y = y_t - y_m
  \]

  \item Compute the Euclidean distance:
  \[
    d = \sqrt{\Delta x^2 + \Delta y^2}
  \]

  \item Compute the target bearing angle:
  \[
    \theta_{\text{target}} = \text{atan2}(\Delta y, \Delta x)
  \]

  \item Compute the forward heading error:
  \[
    e_f = \text{normalize}(\theta_{\text{target}} - \theta_m)
  \]
  where $\text{normalize}(\cdot)$ wraps the angle to $[-\pi, \pi]$.

  \item If \texttt{allow\_reverse} is enabled, also compute the backward heading error:
  \[
    e_b = \text{normalize}(e_f - \pi)
  \]
  and select the direction with the smaller absolute error:
  \[
    \text{direction} = \begin{cases} -1 & \text{if } |e_b| < |e_f| \\ +1 & \text{otherwise} \end{cases}, \quad e_{\text{heading}} = \begin{cases} e_b & \text{if } |e_b| < |e_f| \\ e_f & \text{otherwise} \end{cases}
  \]

  If \texttt{allow\_reverse} is disabled (the default), then $\text{direction} = +1$ and $e_{\text{heading}} = e_f$.
\end{enumerate}

\subsection{Control Decision}

The control output is determined as follows:

\begin{itemize}
  \item \textbf{Catch detection}: If $d < \texttt{catch\_distance}$, the goal succeeds. The turtle stops and the result is returned.

  \item \textbf{Pure rotation}: If $|e_{\text{heading}}| > \texttt{angle\_tolerance}$ (default: 0.1\,rad):
  \[
    v = 0, \quad \omega = \text{clamp}(\texttt{angular\_speed} \cdot \text{sign}(e_{\text{heading}}),\ -\omega_{\max},\ \omega_{\max})
  \]
  The turtle turns in place toward the target.

  \item \textbf{Drive with proportional steering}: Otherwise:
  \[
    v = \text{direction} \cdot \texttt{linear\_speed}, \quad \omega = \text{clamp}(2.0 \cdot e_{\text{heading}},\ -\omega_{\max},\ \omega_{\max})
  \]
  The turtle moves forward (or backward) while applying proportional angular correction to maintain its heading toward the target.
\end{itemize}

\subsection{Design Rationale}

The turn-then-drive strategy was chosen for its simplicity and visual clarity. While a pure proportional-integral-derivative (PID) controller could provide smoother trajectories, the turn-then-drive approach produces a clear, easy-to-understand pursuit behaviour that is well-suited for demonstration purposes. The optional reverse driving feature (\texttt{allow\_reverse}) allows the turtle to take the shorter rotational path by driving backward when the target is behind it, trading visual clarity for efficiency.

\section{Following Control Algorithm}

The follower manager uses a \textbf{proportional control} strategy with a distance threshold. At each control tick (default: 20\,Hz), the following computations are performed for each follower in the chain.

\subsection{Leader Assignment}

For follower $i$ in the chain, its leader is determined by position:

\[
  \text{leader}_i = \begin{cases} \texttt{turtle1} & \text{if } i = 0 \text{ (first follower)} \\ \text{follower}_{i-1} & \text{otherwise (subsequent followers)} \end{cases}
\]

This creates a chain where each follower tracks the one immediately ahead of it, rather than all followers tracking the master turtle directly. This approach produces more natural, snake-like movement.

\subsection{Control Law}

Given the follower's position $(x_f, y_f, \theta_f)$ and the leader's position $(x_l, y_l)$:

\begin{enumerate}
  \item Compute the distance to the leader:
  \[
    d = \sqrt{(x_l - x_f)^2 + (y_l - y_f)^2}
  \]

  \item Compute the heading error:
  \[
    \theta_{\text{target}} = \text{atan2}(y_l - y_f, x_l - x_f), \quad e = \text{normalize}(\theta_{\text{target}} - \theta_f)
  \]

  \item Apply the control decision:
  \begin{itemize}
    \item \textbf{Stop}: If $d < \texttt{follow\_distance}$ (default: 0.8\,m), the follower stops: $v = 0$, $\omega = 0$.
    \item \textbf{Pure rotation}: If $|e| > \texttt{angle\_tolerance}$ (default: 0.15\,rad):
    \[
      \omega = \text{clamp}(\texttt{angular\_speed} \cdot \text{sign}(e),\ -\omega_{\max},\ \omega_{\max})
    \]
    \item \textbf{Proportional drive}: Otherwise:
    \[
      v = \text{clamp}(K_p^{(l)} \cdot (d - d_{\text{follow}}),\ 0,\ v_{\max})
    \]
    \[
      \omega = \text{clamp}(K_p^{(\omega)} \cdot e,\ -\omega_{\max},\ \omega_{\max})
    \]
    where $K_p^{(l)} = 1.0$ and $K_p^{(\omega)} = 4.0$ are the proportional gains.
  \end{itemize}
\end{enumerate}

\subsection{Design Rationale}

The proportional control with distance threshold was chosen to prevent oscillation and collision. Without the distance threshold, followers would attempt to overlap with their leaders, causing erratic behaviour. The proportional gains are tuned to provide responsive following without overshooting. The chain-based leader assignment (each follower tracks its predecessor rather than the master) produces a natural, snake-like formation that is visually appealing and physically plausible.
